% Upper-case    A B C D E F G H I J K L M N O P Q R S T U V W X Y Z
% Lower-case    a b c d e f g h i j k l m n o p q r s t u v w x y z
% Digits        0 1 2 3 4 5 6 7 8 9
% Exclamation   !           Double quote "          Hash (number) #
% Dollar        $           Percent      %          Ampersand     &
% Acute accent  '           Left paren   (          Right paren   )
% Asterisk      *           Plus         +          Comma         ,
% Minus         -           Point        .          Solidus       /
% Colon         :           Semicolon    ;          Less than     <
% Equals        =           Greater than >          Question mark ?
% At            @           Left bracket [          Backslash     \
% Right bracket ]           Circumflex   ^          Underscore    _
% Grave accent  `           Left brace   {          Vertical bar  |
% Right brace   }           Tilde        ~

%
% NSF Biographical Sketch: Dr. Theodore P. Pavlic
%
% Copyright (c) 2013 by Theodore P. Pavlic
% All rights reserved.
%

% \documentclass[svgnames,11pt]{article}
\documentclass[11pt]{article}
\usepackage[
  top    = 0.75in,      % 上边距
  bottom = 0.95in,      % 下边距
  left   = 0.75in,      % 左边距
  right  = 0.75in       % 右边距
]{geometry}
\usepackage[T1]{fontenc}
\usepackage{times}
\usepackage{graphicx,subfig}
\renewcommand\thesection{(\alph{section})}
\renewcommand\thesubsection{(\alph{subsection})}

\usepackage{calc}

\usepackage{lastpage}
%\newcommand{\pageoflastpage}{Page {\thepage} of \pageref*{LastPage}}
%\newcommand{\pageoflastpage}{{\thepage} of \pageref*{LastPage}}
\newcommand{\pageoflastpage}{}

% Redefine the plain pagestyle for the title page
\makeatletter
\let\oldps@plain\ps@plain
\renewcommand{\ps@plain}{\oldps@plain%
\renewcommand{\@evenfoot}{\hfil\pageoflastpage\hfil}%
\renewcommand{\@oddfoot}{\@evenfoot}}
\makeatother

% Use fancy for non-title pages
\usepackage{fancyhdr}
\fancyhead{}
\fancyfoot{}
\cfoot{\pageoflastpage}
\pagestyle{fancy}

\usepackage{xspace}

\usepackage[shortlabels]{enumitem}
\usepackage{graphicx}
\usepackage{xltabular}
\usepackage{makecell}
\usepackage{varioref}
\usepackage[%
        colorlinks=true,urlcolor=,
        pdfpagelabels=true,hypertexnames=true,
        plainpages=false,naturalnames=true,
        ]{hyperref}
%
\usepackage{url}
\newcommand\doilink[1]{\href{http://dx.doi.org/#1}{#1}}
\newcommand\doi[1]{DOI:\doilink{#1}}
\usepackage{cite}

\makeatletter
\newlength{\bibhang}
\setlength{\bibhang}{1em}
\newlength{\bibsep}
 {\@listi \global\bibsep\itemsep \global\advance\bibsep by\parsep}
\newlist{bibsection}{itemize}{3}
\setlist[bibsection]{label=,leftmargin={\bibhang+\widthof{[9]}},%
        itemindent=-\bibhang,
        itemsep=\bibsep,parsep=\z@,partopsep=0pt,
        topsep=0pt}
\newlist{bibenum}{enumerate}{3}
\setlist[bibenum]{label=\textbf{\arabic*.},leftmargin={\bibhang+\widthof{[999]}},%
        itemindent=-\bibhang,
        itemsep=\bibsep,parsep=\z@,partopsep=0pt,
        topsep=0pt}
\let\oldendbibenum\endbibenum
\def\endbibenum{\oldendbibenum\vspace{-.6\baselineskip}}
\let\oldendbibsection\endbibsection
\def\endbibsection{\oldendbibsection\vspace{-.6\baselineskip}}
\makeatother

\title{%
\vspace*{-2\baselineskip}
\noindent
\begin{minipage}[c]{0.73\textwidth} % 左侧占 73%
    \centering
    {\LARGE\textbf{Yiwei Zheng}}\\[0.25em]
    {\normalsize Postdoctoral Researcher}\\[0.8em]
    {\normalsize School of Automation and Intelligent Sensing, Shanghai Jiao Tong University, China}\\
    {\normalsize Date of Birth: Dec. 1, 1996}\\
    {\normalsize e-mail: \href{mailto:ywzheng1201@sjtu.edu.cn}{ywzheng1201@sjtu.edu.cn}}\\
    {\normalsize tel: +86-136-0955-9421}
\end{minipage}%
\hfill
\begin{minipage}[c]{0.23\textwidth} % 右侧占 23%
    \centering
    % 照片大小：与文字高度接近
    \includegraphics[height=4.0cm,keepaspectratio,clip,trim=0 0 0 0]{bio/ywzheng_3.jpg}
\end{minipage}

\vspace{0.6em}
\hrule height 0.6pt % 横线全宽、贯穿整个页面内容宽度
%\vspace{-1.0em}
}
\date{}
\author{}

%\hypersetup{%
%        pdfsubject={NSF Biographical Sketch of Dr.~Theodore P.~Pavlic},
%        pdfauthor={Dr. Theodore P.~Pavlic},
%        pdftitle={NSF Biographical Sketch of Dr.~Theodore P.~Pavlic},
%        pdfkeywords={}
%        }

\newcommand{\ie}{i.e.\xspace}
\newcommand{\eg}{e.g.\xspace}

\begin{document}

\maketitle
\vspace{-4\baselineskip}



\subsection*{Research Interest}

%Multi-agent Path Planning;
%Formal Methods for Dynamic Systems;
%Security of Discrete-Event System.

\begin{bibenum}[itemsep=5pt]
    \item Multi-agent path and mission planning using formal methods.
    \item Safety, security, and opacity of discrete-event systems.
    \item Control \& navigation integration for robotic systems in unknown environments.
\end{bibenum}



\subsection*{Education}
    \begin{xltabular}{\textwidth}{lX}
        2018 -- 2025 & 
        Ph.D., Control Science \& Engineering, Xiamen University, Xiamen, China. \newline 
        (Supervisor: \textbf{Prof. Weiyao Lan}, \textbf{Prof. Xiao Yu}) \\
        
        2014 -- 2018    & 
        B.S., Opto-Electronic Science \& Engineering, Huaqiao University, Quanzhou, China.
    \end{xltabular}


%%
%\subsection*{Appointments}

%\begin{itemize}[label={\quad 9999--9999:},leftmargin=*,itemsep=0pt]

%    \item[2012--present:]
%        \textbf{Postdoctoral Scholar},
%        Arizona State University,
%        Austin, TX

%    \item[2010--2012:]
%        \textbf{Postdoctoral Researcher},
%        The Ohio State University,
%        Columbus, OH

%    \item[2002, 2003:]
%        \textbf{Hardware R\&D Intern},
%        National Instruments,
%        Austin, TX

%    \item[2001:]
%        \textbf{Core Systems Software Developer},
%        IBM Network Storage,
%        Research Triangle Park, NC

%    \item[1997--2001:]
%        \textbf{Information Technology Systems Engineer},
%        CallTech Communications,
%        Columbus, OH

%    \item[1995--1997:]
%        \textbf{Web Developer},
%        MegaLinx Communications,
%        Dublin, OH

%\end{itemize}

\subsection*{Publications}
\begin{bibenum}[itemsep=5pt, label=\textbf{[\arabic*] }]


% for reference
    %\bibliographystyle{IEEEtran}
    %\bibliography{ref.bib}
    %\cite{hou2019model}
    %\cite{zheng2022early}
    %\cite{zheng2022enforcement}
    %\cite{xie2023adaptive}
    %\cite{xie2023data}
    %\cite{zheng2022optimal}
    %\cite{zheng2019hybrid}

    \item \textbf{Y. Zheng}, A. Lai, W. Lan, X. Yu, and R. Su,
    ``Enforcement of opacity for interval weighted discrete-event system by supervisory control,'' 
    \emph{IEEE Transactions on Automatic Control}.
    \doi{10.1109/TAC.2025.3611993}, 2025. (Full Paper)

    \item \textbf{Y. Zheng}, A. Lai, X. Yu, and W. Lan, 
    ``Early-awareness collision avoidance in optimal multi-agent path planning with temporal logic specifications,''
    \emph{IEEE/CAA Journal of Automatica Sinica}, vol. 10, no. 5, pp. 1346-–1348, May 2023.
    % \doi{10.1109/JAS.2022.106043}.
    (Video is available \href{https://www.youtube.com/playlist?list=PLWlDMqOM9D2QprwUcfTXYbPa6IsiOVBEZ}{\textbf{here}}.)

    \item \textbf{Y. Zheng}, Y. Wang, W. Lan, and X. Yu,
    ``Optimal synthesis for partially-observable stochastic systems under opaque temporal logic tasks.'' 
    Submitted.
    (Video is available \href{https://www.youtube.com/playlist?list=PLWlDMqOM9D2Q1lAKcFA_GbEFAeicVcVh2}{\textbf{here}}.)

    \item K. Xie, \textbf{Y. Zheng}, X. Yu, and W. Lan, 
    ``Adaptive optimal output regulation of unknown linear continuous-time systems by dynamic output feedback and value iteration,''
    \emph{Control Engineering Practice}, vol. 141, pp. 105675, Dec. 2023.
    (Video is available \href{https://github.com/trigger1996/Unknown_System_Regulation_ADP}{\textbf{here}}.)

    \item K. Xie, \textbf{Y. Zheng}, W. Lan, and X. Yu, 
    ``Optimal dynamic output feedback control of unknown linear continuous-time system by adaptive dynamic programming,''
    \emph{Automatica}, vol. 163, article 111601, May 2024.
    %(Video is available \href{https://xiaoyu.xmu.edu.cn/demos/UAV_OPFB_IE.mp4}{\textbf{here}}.)

    \item \textbf{Y. Zheng}, W. Lan, and X. Yu, 
    ``Optimal synthesis for stochastic systems
    with information security preservation under temporal logic constraints,''
    in \emph{Proceedings of the 2024 18th International Conference on Control, Automation, Robotics and Vision (ICARCV)},
    Dubai, UAE, Dec. 2024, pp. 918-–923.

    \item \textbf{Y. Zheng}, A. Lai, W. Lan, and X. Yu, 
    ``Optimal path planning with opacity-preserving temporal logic specifications using bipartite synthesizers,'' 
    in \emph{Proceedings of the 62nd IEEE Conference on Decision and Control (IEEE CDC 2023)}, Singapore, Dec. 2023, pp. 7862–-7867.
    
    \item \textbf{Y. Zheng}, Z. Liu, A. Lai, X. Yu, and W. Lan, 
    ``Optimal multi-agent collision-free path planning with temporal logic constraints,'' 
    in \emph{Proceedings of 2021 5th Chinese Conference on Swarm Intelligence and Cooperative Control (CCSICC 2021)}, %\quad Springer, 
    Shenzhen, China, Nov. 2022, pp. 1468–-1477.
    % DOI:\href{http://doi.org/10.1007/978-981-19-3998-3_137}{10.1007/978-981-19-3998-3\_137}.
    
    \item \textbf{Y. Zheng}, Y. Xu, J. Zhang, and D. Luo, 
    ``A hybrid slam method for indoor micro aerial vehicles,''
    in \emph{2019 IEEE 15th International Conference on Control and Automation (ICCA 2019)}, %\quad IEEE
    Edinburgh, Scotland, July 2019, pp. 651–-656.
    % \doi{10.1109/ICCA.2019.8900012}.
    (Video is available \href{https://www.youtube.com/playlist?list=PLWlDMqOM9D2R34O04VglkfJLh-Xs-Wbzt}{\textbf{here}}.)


\end{bibenum}

%\subsection*{Attended Research Projects}

%\begin{bibenum}[itemsep=5pt]
%    \item Project, the National Key R \& D Program of China (Grant No. 2021ZD0112602); 
%    \item General Program, National Natural Science Foundation of China (Grant No. 62173283); 
%    \item General Program, National Natural Science Foundation of China (Grant No. 61673327); 
%    \item Project, Aeronautical Science Foundation of China (Grant No. 20185568005).
%\end{bibenum}

\subsection*{Honors and Awards}

\begin{bibenum}[itemsep=5pt]
    \item Awarded Second Prize Scholarship, Huaqiao University (2014--2015).
    \item Awarded First Prize Scholarship, Huaqiao University (2015--2016, 2016--2017).
\end{bibenum}



\subsection*{Contest Experience}

\begin{bibenum}[itemsep=5pt]
    \item \textbf{Y. Zheng.}, B. Zhao, S. Huang, Z. Huang, and M. Chen, 2022 China Graduate Future Filght Vehicle Innovation Competition, \textbf{3rd Prize}.
    \item \textbf{Y. Zheng.}, B. Zhao, S. Huang, and X. Wang, 2023 China Graduate Future Filght Vehicle Innovation Competition, \textbf{3rd Prize}.
    \item B. Zhao, Z. Huang, J. Guo, X. Chen, \textbf{Y. Zheng.}, W. Lan,  and  X. Yu, 62nd IEEE CDC Cooperative Aerial Robots Inspection Challenge, \textbf{4th Place}.
    \item \textbf{Y. Zheng.}, Z. Xu, and L. Lin, 2017 National Undergraduate Electronics Design Contests, \textbf{2nd Prize in Division}.
    \item \textbf{Y. Zheng.}, Z. Xu, and  L. Lin, 2016 Fujian Undergraduate Electronics Design Contests, \textbf{1st Prize}.
    \item \textbf{Y. Zheng.}, S. Lei, and H. Long, 2015 National Undergraduate Electronics Design Contests, \textbf{2nd Prize in Division}.



\end{bibenum}

\subsection*{Abilities for Unmanned Systems}


My overall abilities for designing, tuning, and operating a unmanned system is listed in detail as follows.

\vspace{3mm}

\begin{figure}[h!]
    \centering
    \includegraphics[width=0.8\linewidth]{pics/Drone_usage.pdf}
    %\caption{\small My overall research abilities and interests, ranging from hardware-level implementation to high-level planning.}
    \label{fig:2}
\end{figure}

\begin{bibenum}[itemsep=5pt]

    \item Design and development of robotic systems (e.g., wheeled robots and unmanned aerial vehicles), encompassing mechanical and electronic hardware as well as mid-level embedded systems.
    
    \item Implementation of control, observation, and planning frameworks derived from theoretical results in formal synthesis and discrete-event systems, validated through physical experiments on robotic platforms.
    


    %\item Small team leading and management.
\end{bibenum}



    %\item \textbf{Conference service:}
    %    (i) Organizer and associate editor for invited session:
    %    ``Correctness by Verification and Design'',
    %    14\textsuperscript{th} IEEE Conference on Intelligent
    %    Transportation Systems, Washington, DC, October
    %    5--7, 2011.
    %    (ii) Organizer for proposed mini-symposium: ``Optimization and
    %    Rationality in Eusocial Insects'', 2013 Society for Mathematical
    %    Biology Annual Meeting and Conference, Tempe, AZ,
    %    June 10--13, 2013
    %    (iii) Co-organizer for technical session: ``Complex Systems of
    %    Social Insects in Research and Education'', 2013 Biomathematics
    %    and Ecology: Education and Research, Arlington, VA,
    %    October 11--13, 2013.

    %\item \textbf{University service:} Chair of committee for the
    %    development of biomimicry and bio-inspired research and education
    %    initiatives at Arizona State University, 2013.

    %\item \textbf{Community service:}
    %    (i) As NSF GK-12 fellow, developed and implemented science
    %    lessons for local inner-city public school, 2006--2007.
    %    (ii) Recognized in \emph{Wired} magazine for OSS work Linux
    %    Virtual Server project, 2000.

    %\item \textbf{Mentoring:}
    %    (i) Founding faculty co-adviser for Interdisciplinary Complexity
    %    Science Student Organization at Arizona State University, 2013.
    %    (ii) Founding mentor of 2002 FIRST high-school robotics team
    %    that received regional honor for inspiration in engineering,
    %    Dublin, OH.
    %    (iii) Founding mentor of 2001 FIRST LEGO League all-girls
    %    middle-school robotics team, Columbus, OH

    %\item \textbf{Teaching:}
    %    (i) Authored hundreds of pages of Creative-Commons-licensed
    %    course material for courses in analog electronics and control
    %    for undergraduate and graduate students, The Ohio State
    %    University, 2007--2009.
    %    (ii) Developed and implemented Electrical and Computer
    %    Engineering undergraduate student design project on
    %    retrofittable vehicle-to-vehicle communications system for
    %    adaptive cruise control in mixed-traffic environments, The Ohio
    %    State University, 2012.
    %    (iii) Mathematical modeling guest lecturer, Research Strategies
    %    for Animal Behavior, Arizona State University, 2013.



\end{document}